\documentclass[10pt,a4paper]{amsart}
\usepackage[margin=19mm]{geometry}
\usepackage{amsmath,amssymb,mathtools}
\usepackage[scr=rsfs,cal=euler]{mathalfa}
\usepackage{xfrac}
\usepackage[unicode,hidelinks,bookmarks=false]{hyperref}
\newcommand{\ie}{i.e.\ }
\newcommand{\cf}{cf.\ }
\newcommand{\resp}{respectively\ }
\newcommand{\Subsection}{\subsection}
\providecommand{\nobreakdash}{\nobreak\hbox{-}}
\setlength{\emergencystretch}{2em}
\multlinegap0pt
\usepackage{mathrsfs}
\usepackage{verbatim}
\usepackage{subfigure}
\usepackage{psfrag}
\usepackage{enumitem}
\newtheorem{theorem}{Theorem}[section]
\newtheorem{lemma}[theorem]{Lemma}
\newtheorem{corollary}[theorem]{Corollary}
\newtheorem{proposition}[theorem]{Proposition}
\newtheorem{conjecture}{Conjecture}
\newtheorem{definition}[theorem]{Definition}
\newtheorem{example}[theorem]{Example}
\newcommand{\Z}{\mathbb{Z}}
\newcommand{\R}{\mathbb{R}}
\newcommand{\C}{\mathbb{C}}
\newcommand{\tatop}[2]{\genfrac{}{}{0pt}{10}{#1}{#2}}
\newcommand{\Res}{\textnormal{Res}}
\newcommand{\Root}{\textnormal{Root}}
\begin{document}
\title[On mappings with Jacobian one]{On mappings with Jacobian one}
\author{Zbigniew Jelonek}
\date{}
\maketitle
\noindent\emph{Source and licence.} On mappings with Jacobian one. Version arXiv:2607.20597v1 (CC BY 4.0). This material is distributed under the Creative Commons Attribution 4.0 International licence, http://creativecommons.org/licenses/by/4.0/, as recorded in the arXiv OAI licence metadata for this exact version and shown on the abstract page https://arxiv.org/abs/2607.20597v1. Full rights notice: licenses/arxiv-2607.20597-CC-BY-4.0.txt.\par\smallskip
\noindent\textit{Editorial scope.} Counted unit: the original \texttt{theorem} environment \texttt{-} at source lines 108--109 together with its complete proof at lines 111--113; the delivered document keeps source lines 76--114 so that every referenced label and every citation resolves inside it. Native editable source is the paper's own concatenated TeX; the AMS wrapper is editorial formatting only. No publisher class, upstream script or unrelated artwork is redistributed.
\section{Main Result}

\begin{lemma}
The set $A(n, d)$ of polynomial automorphisms $F : \Bbb C^n \to \Bbb C^n$ of degree at most  $d$ and with $Jac(F ) = 1$ is Zariski closed.
\end{lemma}

\begin{proof}
It is easy to see that the set $A(n, d)$ is a constructible subset of $\Omega(n,d)\cong\Bbb C^M$ (where  $\Omega(n,d)$ is a space of all  mappings $F:\Bbb C^n \to \Bbb C^n$ of degree at most $d$).  Indeed, consider the algebraic set 
$$S=\{ (F,G)\in \Omega(n,d)\times \Omega(n, d^{n-1}) ; F\circ G=identity\}.$$ Then the projection $\pi(S)$ of
$S$ on  $\Omega(n,d)$ is a constructible set by Chevaley theorem. But $A(n,d)=\pi(S) \cap \{ F : Jac(F)=1\},$ hence it is a constructible set.



Hence it is enough to prove that the set $A(n, d)$ is closed in the euclidian topology. Let $f_k \in A(n, d)$ and $f_k \to F.$ 
Hence $F = (F_1, ..., F_n)$ is a polynomial mapping with $Jac(F ) = 1.$ We can assume that $f_n (0) = 0.$ Indeed, $f_k (0) = b_k$ and $b_k \to b = F (0).$

 Now if we show that $f_k - b_k \to G \in A(n, d)$,  then also $f_k \to G + b.$ Now let $f_k(1, ..., 1) = (w_{k,1} , ..., w_{k,n}) = w_k \in \Bbb C^n.$
  Hence $w_k \to w = F (1, ..., 1).$  If $g_k = f^{ -1}_k,$ then $g_k(w_k) = (1, ..., 1).$  Note that the family of all $w_k$ is bounded.
  Now let $g_k = (Q_{k,1} , ..., Q_{k,n}).$ It is well known that the algebraic degree of $g_k$ is uniformly bounded by $d^{n-1}.$
  
   For a polynomial $Q = \sum a_\alpha x^\alpha\in  \Bbb C[x_1, ..., x_n]$ let $||Q|| =max |a_\alpha|.$ Let us denote $\tilde{Q} = Q/||Q||.$
    Since $Q_{k,j} (w_k) \to 1$ and $w_k$ is bounded, we have that there is an $\epsilon > 0$, such that $||Q_{k,j} || > \epsilon$ for every $k, j.$
    
    Now we have $Q_{k,1} (f_{k,1} , ..., f_{k,n} ) = x_1$. Consequently $\tilde{Q}_{k,1} (f_{k,1} , ..., f_{k,n} ) = a(k) x_1$, where $0 <a(k) < 1/\epsilon$. 
    If we take a subsequence of $\tilde{Q}_{k,1}$ we can assume that it converges to a polynomial $Q_1$ and one of its (fixed) coefficients has norm $1$. 
    Moreover we can assume that $a(k)$ also converges to $a$. Let us note that $a  \not= 0,$ because otherwise $Q_1(F_1, ..., F_n) = 0$ and $Q_1$ is a non-constant polynomial. 
    
    This contradicts the fact that $F_1, ..., F_n$ are algebraically independent, as components of a mapping with Jacobian equal to $1.$
     Hence finally $(1/a)Q_1(F_1, ..., F_n) = x_1,$ i.e., there is a polynomial $Q'_1$ such that $Q'_1(F_1, ..., F_n) = x_1.$
      In the same way there are polynomials $Q'_2, ..., Q'_n$ such that $Q'_j (F_1, ..., F_n) = x_j$ for $j = 2, ...n.$ This means that F is an automorphism.
\end{proof}


\begin{theorem} Let $X(n,d)$ denotes the space (which is an affine algebraic variety) of polynomial mappings $F$ with $Jac(F)=1$ and of degree bounded by $d.$ Let $Y$ be an irreducible component of this set. Then either $Y$ consists of polynomial automorphisms, or a generic element of $Y$ is a counterexample to Jacobian Conjecture. Moreover, $\dim Y \ge n^2-1.$
\end{theorem}

\begin{proof} Let $Y$ be an irreducible component. Then either $Y\subset A(n,d)$ or it contains a counterexample to the Jacobian Conjecture. In the latter case since $A(n,d)$ is closed it means that $A(n,d)\cap Y$ is a nowhere dense in $Y.$ Now denote by $\mathcal L$ the set of linear automorphisms with the Jacobian $1.$ Take $F\in Y$, which is not contaned in any different component of $X(n,d).$
Then $F\circ \mathcal{L}$ is an irreducible set which is contained in $X(n,d).$ Since it has a common point which belongs to $Y$ and does not belong to any other component it means that $F\circ \mathcal{L}\subset Y.$ But $\dim F\circ \mathcal{L}=n^2-1.$
\end{proof}


\end{document}
